\chapter{Supervision and Guidance Strategies for SVG Generation}\label{ch:supervision}

While representation defines how vector graphics are encoded during 
generation, supervision determines how those representations are shaped. 
In SVG generation, supervision refers to the signals that guide learning 
or optimization so that produced vector graphics align with a target 
objective. These signals may come from paired datasets, reconstruction 
losses, semantic guidance models, or structural priors introduced during 
initialization.

An important distinction is that supervision in this domain extends 
beyond dataset labels. Many methods do not rely on ground-truth SVG 
targets at all. Instead, they optimize vector representations using 
raster reconstruction, semantic similarity, or pretrained generative 
models as guidance. Supervision is therefore a core architectural 
decision: it determines which aspects of the graphic are emphasized 
during generation and shapes how structure emerges.

Supervision also interacts closely with representation. Because 
gradients propagate through the chosen representation, the same 
supervision signal can produce different geometric outcomes depending 
on how the vector graphic is encoded. Image-based losses applied to 
explicit primitives may introduce additional paths to match appearance, 
whereas the same losses applied to implicit neural representations may 
reshape continuous fields. This interaction highlights that supervision 
and representation need to be analyzed together.

A further distinction concerns when supervision is applied. Some models 
learn reusable generation behavior from datasets at training time. 
Others apply supervision directly during inference, refining a vector 
representation for each individual input. This difference has 
significant consequences for speed, structure, reproducibility, and 
practical usability, and is examined in detail in Section~\ref{sec:dataset_vs_instance}.

The following sections examine the main supervision families used 
across the selected models, focusing on how different guidance 
strategies influence structure, editability, visual quality, and 
practical usability.


\section{Reconstruction-Based Supervision}

Reconstruction-based supervision guides generation by comparing a 
rendered SVG output with a target raster image. The objective is to 
minimize the difference between prediction and reference using loss 
functions defined in image or feature space. Common choices include 
pixel-wise L1 or L2 losses and perceptual losses computed from 
pretrained convolutional networks~\cite{johnsonPerceptualLossesRealTime2016}. 
This paradigm is particularly common in dataset-trained models and 
image-to-SVG approaches.

One advantage of reconstruction-based supervision is stability. 
Because the target is known, gradients provide direct information 
about how geometry should change. Training behavior is predictable, 
and when datasets contain clean SVG examples, models can learn 
typical path counts, grouping patterns, and structural conventions. 
The dataset itself acts as the primary supervision signal: models 
such as DeepSVG~\cite{carlierDeepSVGHierarchicalGenerative2020}, 
DeepIcon~\cite{bingDeepIconHierarchicalNetwork2024}, and Im2Vec~\cite{reddyIm2VecSynthesizingVector2021} are trained to 
reconstruct vector graphics from data, and the reconstruction 
objective is what drives learning. As discussed in 
Chapter~\ref{ch:Landscape}, these models do not apply supervision 
at inference time, once trained, generation occurs through a 
forward pass.

However, reconstruction-based supervision has limits. Output quality 
depends heavily on dataset distribution: models trained on simplified 
icons may generalize poorly to complex scenes, and datasets with 
inconsistent SVG structure can transfer those inconsistencies to 
generated outputs. More importantly, reconstruction losses measure 
visual similarity rather than structural correctness. A model can 
achieve low reconstruction error while producing redundant paths or 
weak hierarchy, particularly when pixel-level matching dominates the 
objective. Without explicit structural constraints, reconstruction 
supervision does not guarantee editability or geometric clarity.


\section{Semantic Guidance with Vision–Language Models}

Many SVG generation tasks operate without paired SVG references. 
In text-to-SVG settings, the only available signal may be a textual 
description of the desired graphic. Semantic guidance addresses this 
by using pretrained vision–language models as evaluators that provide 
feedback about whether the rendered SVG aligns with a prompt.

As described in Chapter~\ref{ch:Background}, CLIP embeds images and 
text into a shared representation space~\cite{radfordLearningTransferableVisual2021}. 
In SVG generation pipelines, the rendered vector graphic is passed 
through the image encoder and compared against the encoded text 
prompt. Because gradients can propagate back through the 
differentiable rasterizer, vector parameters can be adjusted to 
improve semantic alignment. This enables open-ended generation 
without requiring paired SVG datasets.

Semantic guidance offers flexibility across diverse prompts and 
captures high-level visual semantics that are difficult to encode 
through handcrafted losses. At the same time, it remains indirect: 
the signal measures alignment in embedding space rather than 
geometric correctness. Optimization may therefore introduce 
additional primitives or unstable geometry to satisfy semantic 
objectives, without any guarantee of clean structure or editability.

In practice, semantic guidance is rarely used in isolation. It 
typically acts as an intermediate step or is combined with 
diffusion-based supervision, as in VectorFusion~\cite{jainVectorFusionTexttoSVGAbstracting2022} 
and early SVGDreamer configurations. This reflects a broader 
pattern: CLIP guidance is most useful for conceptual alignment, 
but stronger perceptual signals are usually needed to achieve 
visually convincing results.


\section{Diffusion-Based Supervision}

Diffusion-based supervision leverages pretrained diffusion models 
as perceptual priors rather than relying on explicit targets or 
embedding similarity alone. These models capture large-scale image 
distributions and provide gradients that encourage rendered SVGs 
to move toward visually plausible solutions conditioned on a 
text prompt.

The most common implementation is Score Distillation Sampling (SDS), 
introduced by DreamFusion~\cite{pooleDreamFusionTextto3DUsing2022}. 
SDS removes the need for ground-truth images by treating a frozen 
diffusion model as the optimization signal --- it evaluates the 
rendered SVG and provides gradients that push vector parameters 
toward outputs the diffusion model considers plausible for the 
given prompt. Variational Score Distillation (VSD), used in 
SVGDreamer~\cite{xingSVGDreamerTextGuided2024}, refines this 
further by reducing the oversaturation artifacts that plain SDS 
tends to produce.

Several models in this thesis use diffusion-based supervision 
as their primary guidance signal, including VectorFusion, 
SVGDreamer~\cite{xingSVGDreamerTextGuided2024}, SVGDreamer++~\cite{xingSVGDreamerAdvancingEditability2025}, 
NeuralSVG~\cite{polaczekNeuralSVGImplicitRepresentation2025}, NIVeL~\cite{thamizharasanNIVeLNeuralImplicit2024}, 
and T2V-NPR~\cite{zhangTexttoVectorGenerationNeural2024}.

The main advantage is expressive power. Because diffusion models 
encode rich visual knowledge, they can guide generation toward 
complex compositions and stylistic diversity that reconstruction 
losses alone cannot achieve. However, diffusion gradients operate 
in pixel space and primarily emphasize perceptual realism rather 
than geometric simplicity. Without additional structural 
constraints, optimization tends to increase primitive count to 
satisfy visual objectives, which directly affects editability. 
Inference cost is also substantially higher than for forward-decoding 
approaches, since each generation requires repeated rendering and 
evaluation cycles.

Despite these trade-offs, diffusion-based supervision has become 
the dominant strategy for open-ended SVG generation because it 
removes the dependency on paired datasets while providing strong 
perceptual guidance. Most recent models combine it with 
representation innovations or structural priors to mitigate 
geometric instability, which is why supervision cannot be 
analyzed independently from the other design choices examined 
in this thesis.


\section{Dataset Supervision vs.\ Per-Instance Supervision}
\label{sec:dataset_vs_instance}

A fundamental distinction across SVG generation methods is when 
supervision is applied. Dataset-supervised approaches learn 
generation behavior from collections of SVG graphics during 
training. Per-instance supervision applies guidance during 
generation itself, refining a vector representation for each 
individual input or prompt.

For dataset-supervised models, the dataset is the supervision 
signal. Models such as DeepSVG~\cite{carlierDeepSVGHierarchicalGenerative2020}, DeepIcon~\cite{bingDeepIconHierarchicalNetwork2024}, and Im2Vec~\cite{reddyIm2VecSynthesizingVector2021} are trained 
on reconstruction objectives against SVG data, and once trained, 
generation requires no further optimization. LayerTracer~\cite{songLayerTracerCognitiveAlignedLayered2025} takes 
a specific variant of this approach: it is trained on a 
synthetically generated serpentine dataset, where vector 
graphics are constructed in a deliberate layer order that 
mimics human drawing behavior~\cite{songLayerTracerCognitiveAlignedLayered2025}. 
The dataset structure itself encodes the supervision prior, 
teaching the model to generate layered, sequential outputs 
that reflect design conventions rather than arbitrary path 
orderings.

Per-instance supervision works differently. Models such as 
LIVE~\cite{maLayerwiseImageVectorization2022}, 
SAMVG~\cite{zhuSAMVGMultistageImage2023}, VectorFusion~\cite{jainVectorFusionTexttoSVGAbstracting2022}, 
SVGDreamer~\cite{liSVGDreamerTexttoSVG2023}, NeuralSVG~\cite{liNeuralSVGTexttoSVG2023}, NIVeL~\cite{zhaoNIVeLTexttoVector2023}, and T2V-NPR~\cite{liT2V-NPRTexttoVector2023} apply supervision 
directly during generation. No reusable generation model is 
trained; instead, vector parameters are optimized from scratch 
for each input using reconstruction losses, semantic alignment, 
or diffusion priors. This is computationally expensive but 
allows structure to emerge dynamically in response to each 
specific prompt or image.

Some models combine both paradigms. SVGFusion~\cite{xingSVGFusionScalableTexttoSVG2025} learns 
latent representations from SVG datasets during training, then 
applies optimization-based refinement at inference time. This 
hybrid approach attempts to balance the structural consistency 
of dataset supervision with the flexibility of per-instance 
optimization. The same pattern appears in T2V-NPR~\cite{zhangTexttoVectorGenerationNeural2024}, which 
learns path-level representations from data before using 
diffusion guidance during generation.

This distinction has direct practical implications. 
Dataset-supervised models produce consistent outputs and 
support fast inference, but their structure is constrained by 
what the training data contains. Per-instance supervision 
offers greater adaptability but introduces variability due 
to initialization choices and stochastic guidance, and 
makes systematic evaluation across methods more difficult. 
Hybrid approaches mitigate some of these trade-offs but 
increase overall system complexity.


\section{Structural Priors and Initialization as Supervision}

Beyond explicit loss functions, many SVG generation methods 
rely on structural priors and initialization strategies that 
act as additional forms of supervision. These mechanisms are 
not always described as supervision, but they strongly 
influence how structure emerges in the resulting SVG.

Initialization determines the starting point of optimization. 
Random initialization provides flexibility but often leads to 
unstable optimization and fragmented paths. Structured 
initialization introduces prior knowledge about object 
boundaries and layout, reducing the need for excessive 
refinement. SAMVG~\cite{zhuSAMVGMultistageImage2023} illustrates this by using segmentation 
outputs to define an initial set of semantically meaningful 
regions before vector optimization begins. LIVE~\cite{maLayerwiseImageVectorization2022} and 
SVGDreamer++~\cite{liSVGDreamerTexttoSVG2023} use layered initialization strategies that 
progressively refine structure across stages. In these cases, 
the initialization itself constrains geometry before any 
loss is applied.

Structural priors can also come from learned representations. 
SVGFusion~\cite{xingSVGFusionScalableTexttoSVG2025} and 
T2V-NPR~\cite{liT2V-NPRTexttoVector2023} both use latent starting points shaped 
by prior training, so optimization begins from a representation 
that already encodes typical path structures rather than from 
scratch. This often improves convergence speed and reduces 
structural noise compared to random initialization.

The interaction with diffusion supervision is particularly 
important here. Because diffusion gradients operate in pixel 
space and do not enforce hierarchy, initialization plays a 
critical role in determining whether structure remains 
interpretable. When meaningful priors are absent, 
diffusion-driven optimization tends to increase primitive 
count to match visual detail. When structured initialization 
is present, the same supervision signal can produce more 
coherent geometry with fewer paths. This is one reason why 
models with similar supervision strategies can produce very 
different structural outcomes depending on how generation 
is initialized.

From a broader perspective, treating initialization and 
structural constraints as forms of supervision helps explain 
observations that loss functions alone cannot account for. 
Methods that incorporate segmentation, layering, or learned 
latent starting points often generate more interpretable SVG 
outputs despite relying on similar perceptual objectives as 
less structured alternatives.


\section{Summary}

This chapter analyzed supervision as a key dimension shaping 
SVG generation methods. While representation determines how 
vector graphics are encoded, supervision determines how these 
representations evolve toward desired outputs.

Reconstruction-based supervision provides explicit targets and 
stable optimization, but its quality depends on dataset 
distribution and it does not guarantee structural correctness. 
Semantic guidance enables generation without paired data, but 
its indirect nature can lead to unstable geometry. 
Diffusion-based supervision provides strong perceptual priors 
that improve visual expressiveness significantly, at the cost 
of increased computational load and structural complexity.

The distinction between dataset supervision and per-instance 
supervision clarifies how methods allocate computational effort. 
Dataset-trained models use the training data itself as the 
supervision signal and favor efficiency and consistency. 
Per-instance optimization offers flexibility and expressive 
power at the cost of runtime and reproducibility. Hybrid 
approaches such as SVGFusion~\cite{xingSVGFusionScalableTexttoSVG2025} 
and T2V-NPR~\cite{liT2V-NPRTexttoVector2023} combine both, 
reflecting an ongoing attempt to balance speed with 
adaptability. LayerTracer~\cite{maLayerwiseImageVectorization2022} 
represents a particularly direct example of how dataset design can encode structural priors: 
its serpentine training data teaches the model to produce 
layered outputs without requiring any explicit structural 
loss.

Structural priors and initialization strategies extend the 
definition of supervision beyond explicit loss functions. 
They shape the space in which optimization occurs and 
strongly influence geometry quality, often more than the 
choice of loss function itself.

Across all supervision strategies, a recurring pattern 
emerges: supervision tends to prioritize visual objectives, 
while structural properties emerge indirectly. This explains 
why models with strong perceptual guidance can produce 
visually convincing SVGs that are difficult to edit. When 
analyzed together with representation, supervision provides 
a framework for understanding why different model 
architectures produce distinct structural outcomes. The next 
chapter examines how these properties can be measured and 
compared through evaluation.